\chapter{Installation}

\section{Requirements}

\hayashi{} is distributed as a static binary for Linux, macOS, and
Windows. There are no runtime dependencies: the \hay{} executable is
self-contained.

\begin{itemize}
  \item Linux x86-64 or ARM64
  \item macOS 12+ (Intel or Apple Silicon)
  \item Windows 10+ (x86-64)
\end{itemize}

\section{Via Cargo (recommended)}

If Rust is installed:

\begin{lstlisting}[style=inline, language=sh]
cargo install hayashi-lang
\end{lstlisting}

The \texttt{hay} binary will be installed in \texttt{\$HOME/.cargo/bin}.
Make sure that directory is on your \texttt{PATH}.

\begin{lstlisting}[style=inline, language=sh]
hay --version
\end{lstlisting}

\section{Pre-compiled binaries}

Available on the repository releases page:
\url{https://github.com/sheep-farm/hayashi/releases}

Download the archive for your platform, extract it, and place the
\texttt{hay} binary in any directory on your \texttt{PATH}.

\section{Building from source}

\begin{lstlisting}[style=inline, language=sh]
git clone https://github.com/sheep-farm/hayashi
cd hayashi
cargo build --release
# binary at target/release/hay
\end{lstlisting}

\section{Verifying the installation}

\begin{lstlisting}[style=inline, language=sh]
hay --version
hay --help
\end{lstlisting}

To run a script:

\begin{lstlisting}[style=inline, language=sh]
hay my_script.hay
\end{lstlisting}

To enter the interactive REPL:

\begin{lstlisting}[style=inline, language=sh]
hay
\end{lstlisting}

\begin{nota}
  On Arch Linux with Omarchy/Hyprland, \texttt{\$HOME/.cargo/bin} is
  usually on the \texttt{PATH} if the shell was configured with
  \texttt{rustup}. Otherwise, add to \texttt{.bashrc} or \texttt{.zshrc}:
  \texttt{export PATH="\$HOME/.cargo/bin:\$PATH"}
\end{nota}
